\chapter{Proving the correspondence}
\label{chap:proof}

\section{The theorem and its assumptions}
% State both the expression and matching equivalences, including final-result
% conditions, L = update_locs St, and the shared FreshVarGen parameter.
% Display the proved relations: big-step evaluation iff balanced trees, and
% balanced trees imply machine executions. Do not label the latter an equivalence.

\section{From big-step derivations to balanced trees}
% Explain mutual induction on the big-step derivations. Show how a representative
% rule supplies the transitions and subtrees of its balanced tree constructor.

\section{From balanced trees to big-step derivations}
% Explain the reverse mutual induction and why its matching statement must
% generalize the final argument stack. Use the following cases to show where
% the statement and the tree structure matter.

\subsection{Variable lookup and heap update}
% Reconstruct EvalVar from its subtree and the update transition.
% Track the black-hole set through KUpdate and identify the tree side condition
% used to exclude an already active variable.

\subsection{Matching results and alternatives}
% Explain the distinct treatment of return and failure, including the empty
% final argument stack required for failure. Show how a failed left alternative
% passes its updated heap and restores the saved arguments for the right branch.

\section{From balanced trees to machine executions}
% Explain how trees flatten to the reflexive-transitive closure of step.
% Derive the big-step-to-machine result and its empty-heap, empty-stack instance.
% The initial iff theorem concerns balanced trees, not arbitrary machine traces.

\section{The remaining trace-to-tree direction}
% State the missing decomposition result for ordinary balanced machine traces.
% Compare the checked statements with Theorems 4.1 and 4.5 and Corollary 4.6 of
% the source paper. Separate the proved correspondence from the work needed to
% recover the full machine-to-big-step claim and justify reachability invariants.
